\section{La pirámide de datos: datos reales, de simulación y de internet}

La mayor diferencia entre los datos de robótica y los de los LLM es el anclaje físico. Un LLM puede obtener de internet cantidades enormes de texto, código y material multimodal; un robot necesita acciones, estados, entorno y retroalimentación, y recolectarlos en el mundo real es costoso y lento. Por eso Xie Chen (谢晨) enfatiza que los datos de robótica se parecen más a una pirámide: en la cima están los datos reales recolectados en el propio dispositivo, en el nivel intermedio los datos de simulación o sintéticos, y en la base los datos de internet y de personas.

\begin{figure}[H]
\centering
\includegraphics[width=0.94\textwidth]{data-pyramid.png}
\caption{La pirámide de datos de robótica: datos reales, datos de simulación y datos de internet o de personas.}
\end{figure}

\begin{knowledgebox}{Lectura de la figura: qué resuelve cada uno de los tres niveles de datos}
Los datos reales de la cima son los más cercanos al despliegue, pero son escasos y costosos; los datos de simulación del nivel intermedio son controlables y escalables, adecuados para generar escenarios de cola larga; los datos de internet o de personas de la base tienen la mayor escala y el menor costo, pero su anclaje físico es débil. Una estrategia de datos para robótica tiene que combinar los tres niveles.
\end{knowledgebox}

\subsection{Por qué los datos de robótica son un desierto}

Los robots no son como la conducción autónoma, donde millones de vehículos envían todos los días datos reales de las calles. Los robots de propósito general todavía no se han desplegado a gran escala, la teleoperación es costosa, los escenarios de las tareas están dispersos y los cuerpos de los robots difieren mucho entre sí. Sin escala de despliegue en los dispositivos no surge de forma natural un motor de datos como el de Tesla. Por eso, para la robótica, la simulación y los datos sintéticos son una condición previa, no solo un complemento.

\begin{warningbox}{La simulación no es un juguete, pero tampoco una panacea}
La simulación permite generar escenarios y evaluaciones a escala, pero existe la sim-to-real gap. Los datos reales son escasos, la simulación es imprecisa y los datos de internet tienen un anclaje débil: los tres deben combinarse, no sustituirse entre sí.
\end{warningbox}

\subsection{Resumen de la sección}

La contradicción central de los datos de robótica es esta: los datos más reales son los más escasos, y los más abundantes son los menos reales. La pirámide de datos es el marco básico para resolver esta contradicción.

